\mychapter{I. INTRODUCTION}{I. INTRODUCTION}

\section{Presentation of Internship}

\hs This chapter establishes the context of the thesis. Before turning to the technical
details of models, pipelines, and results, it is useful to explain where the project
originated, why it is relevant, and how it was carried out. The work presented here was
completed during a final year research internship at the Research and Data Analytics
Laboratory (ReDA Lab) of the Institute of Technology of Cambodia (ITC). It represents
the concluding component of the applied mathematics and data science program and brings
together the theoretical and practical knowledge developed throughout the degree.

\hs A final year internship plays an important role in the transition from academic
study to professional practice. It is the stage at which classroom theory is tested
against a real problem that has no predetermined solution. Over a period of four months,
the central aim was to translate concepts from artificial intelligence and natural
language processing into a functional system capable of producing useful results. The
project was conducted under the supervision of mentors from the laboratory, who provided
technical guidance while allowing the independence necessary for genuine research. This
report documents that effort, including both the design decisions that proved effective
and the difficulties encountered along the way.

\hs The internship also served a broader purpose. It provided an opportunity to work in
the manner expected of a practising engineer rather than a student. This involved
deploying software on a real server, handling imperfect data, reviewing research
literature and selecting which methods to adopt, and critically evaluating whether the
resulting system performed well in practice rather than only in appearance. Competencies
of this kind are difficult to acquire through coursework alone and form an essential part
of professional preparation.

\subsection{Objective of Internship}

\hs The objectives of the internship were defined from two perspectives: the academic
requirements of the degree and the personal development goals of the intern. From an
academic standpoint, the internship was intended to fulfil the final requirement of the
applied mathematics and data science program through the completion of a substantial,
independent research project. This required applying knowledge accumulated over the
course of the degree to a problem of practical significance, and presenting the outcome
in a form suitable for academic evaluation.

\hs From a professional standpoint, the internship aimed to develop the competencies
expected of an engineer working in the field of artificial intelligence. The specific
objectives in this regard were to gain hands-on experience in designing and building a
complete AI system rather than an isolated component, to become familiar with modern
tools and techniques in natural language processing and large language models, and to
practise the full development cycle, from initial research and data preparation through
to implementation, evaluation, and deployment.

\hs A further objective was to experience the working environment of a research
laboratory. This included collaborating with supervisors and colleagues, communicating
technical ideas clearly, managing time across overlapping tasks, and responding
constructively to feedback. Taken together, these objectives were intended to ensure
that the internship contributed not only a completed thesis, but also the practical
maturity and professional skills necessary for a career in data science and artificial
intelligence.


\subsection{Duration of Internship}

\begin{table}[H]
\centering
\begin{tabular}{|c|c|}
\hline
\cellcolor{blue!25}{\textbf{Start Date}} & \cellcolor{blue!25}{\textbf{1-March-2026}} \\ \hline
Duration      & 4 months \\ \hline
Location      & Institute of Technology of Cambodia (ITC) \\ \hline
Laboratory    & Research and Data Analytics Laboratory (ReDA Lab) \\ \hline
Position      & AI Engineer \\ \hline
Advisor       & Mr. LONG Pakrigna \\ \hline
Supervisor    & Dr. PHOUK Sokhey \\ \hline
End Date      & 30-May-2026 \\ \hline
\end{tabular}
\caption{Duration of Research}
\label{table:research_duration}
\end{table}

% ─────────────────────────────────────────────────────────────────────
\section{Presentation of Organization}

\hs The Institute of Technology of Cambodia, commonly referred to as ITC, is the
country's leading public engineering university. It was established to provide rigorous
education in science, technology, engineering, and mathematics, and over the years it
has become one of the most respected institutions in these fields in Cambodia. Its
graduates contribute to industry, government, and academic research, both within the
country and internationally.

\hs The institute is led by \textbf{H.E. Prof. Dr. PO Kimtho}, Director of ITC. Under
his leadership, the institution continues to strengthen its academic standards and to
ensure that graduates are well prepared to support national development. This emphasis
extends beyond theoretical instruction toward solving problems of genuine relevance to
Cambodia, an orientation that aligns closely with the aims of the present project.

\hs Within ITC, this work falls under the \textbf{Department of Applied Mathematics and
Statistics}, headed by \textbf{H.E. Dr. LIN Mongkul Sery}. The department provides
training in mathematical modelling, statistics, data science, and artificial
intelligence. It offers a suitable foundation for a project of this nature, since
building an AI system requires both a solid mathematical basis and practical software
engineering skills, and the department's curriculum addresses both.

\subsection{Research and Data Analytics Laboratory (ReDA Lab)}

\hs The research was conducted within the Research and Data Analytics Laboratory (ReDA
Lab), led by \textbf{Dr. PHOUK Sokhey}. The laboratory is a research unit at ITC focused
on applied data science, machine learning, and artificial intelligence. A defining
characteristic of the laboratory is its orientation toward practical impact: rather than
pursuing research purely for publication, it seeks to develop AI solutions for problems
that directly affect people in Cambodia. Much of its activity concerns natural language
processing, predictive analytics, and intelligent systems, which places the present
thesis firmly within its area of expertise.

\hs Conducting the project within ReDA Lab provided access to guidance, computing
resources, and a community of researchers engaged with related problems. This research
environment proved valuable throughout the work. The availability of experienced
colleagues to consult when difficulties arose, and to question assumptions before they
were carried too far, contributed meaningfully to the direction and quality of the final
system.

\subsection{Vision}

\hs To become a leading research laboratory in Southeast Asia for applied data science
and artificial intelligence, producing innovative solutions that address the challenges
of Cambodian society and contribute to national digital transformation.

\subsection{Mission}

\hs \textbf{ReDA Lab} is committed to advancing applied research in data science and
artificial intelligence for social impact. The laboratory's mission is to develop
practical AI solutions tailored to Cambodia's needs, to train the next generation of
data scientists and AI researchers, and to foster academic and industry collaboration
that accelerates the adoption of intelligent technologies across key sectors including
health, education, agriculture, and law.

\subsection{Logo}

%\imgplaceholder[0.4\textwidth]{Institute of Technology of Cambodia Logo}{itc_logo}
\fg{0.3\textwidth}{Institute of Technology of Cambodia Logo}{images/schools/itc.png}
\fg{0.4\textwidth}{Research and Data Analytics Laboratory (ReDA Lab)}{images/schools/ReDA-Lab.png}

\subsection{Organization Chart}

%\imgplaceholder[0.8\textwidth]{ReDA Lab Organization Chart}{reda_org_chart}

\fg{0.6\textwidth}{ReDA Lab Organization Chart}{images/schools/ReDA-org-chart}

\subsection{Address and Contact}
\begin{itemize}
    \item Location: Russian Confederation Blvd, Khan Toul Kork, Phnom Penh, Cambodia
    \item Website: \url{https://www.itc.edu.kh}
    \item Email: \url{info@itc.edu.kh}
\end{itemize}

% ─────────────────────────────────────────────────────────────────────
\section{Background}

\hs Access to legal information is widely regarded as a fundamental right. In Cambodia,
however, this right remains difficult to exercise in practice. The laws exist and are
publicly recorded, yet reading them is far from straightforward. Cambodian legal text is
written in a formal register of Khmer that differs considerably from everyday speech.
The sentences are long, the terminology is specialised, and even well-educated readers
often find it difficult to determine the precise meaning of a given article.

\hs For most people, the conventional solution is to consult a lawyer. This option is
expensive, and for a worker or a small business owner seeking to understand a basic
question about their rights, it is frequently impractical. As a result, many citizens
remain unaware of provisions that directly affect their daily lives, such as the notice
period an employer must provide before termination, or the rate at which overtime should
be paid. The information is technically available to the public, yet effectively beyond
their reach.

\hs In parallel, a significant shift has occurred in the field of artificial
intelligence. Large Language Models (LLMs) have become remarkably capable of reading
text and answering questions about it. Used in isolation, however, these models have a
notable weakness: they can generate information that is incorrect while presenting it
with apparent confidence. For a domain as sensitive as the law, this behaviour poses a
serious risk. Retrieval-Augmented Generation (RAG) was developed to address precisely
this issue. The principle is simple but effective. Rather than allowing the model to
answer from memory, the system first searches a trusted collection of documents,
retrieves the relevant passages, and supplies them to the model as context. The model
then formulates its answer on the basis of what is actually written, producing a
response that can be verified and traced back to its source.

\hs A further difficulty is that most of this progress has taken place in
high-resource languages such as English, French, and Chinese. Khmer, by contrast, is
considered a low-resource language. Fewer tools and datasets are available, and existing
AI models have been exposed to comparatively little Khmer text during training. This
scarcity complicates every stage of the process, from interpreting a question correctly
to matching it with the appropriate legal provision. Consequently, systems designed for
high-resource languages cannot simply be transferred to Khmer without substantial
adaptation.

\hs The present research addresses this gap directly. It proposes an agentic RAG system
designed specifically for Cambodian law, combining two complementary forms of search
with a multi-step reasoning pipeline, so that a question expressed in natural Khmer can
return an accurate answer supported by a citation to the exact article on which it
relies.

% ─────────────────────────────────────────────────────────────────────
\section{Problem Statement}

\hs At present, a Cambodian citizen wishing to consult the law has few practical means
of doing so. Searching through lengthy legal documents by hand is slow, confusing, and
prone to error. No accessible digital tool currently exists that allows a user to pose a
question in natural Khmer and receive a clear, accurate answer accompanied by a reference
to its source.

\hs Developing such a tool is more challenging than it initially appears, for several
reasons. The first is the gap between the way users express questions and the way the
law is written. A query may be phrased informally, whereas the corresponding legal
provision uses formal terminology that shares few words with the question. A simple
keyword search fails to bridge this gap, and even semantic search performs unevenly
because Khmer is a low-resource language. The second difficulty is the absence of an
established Khmer legal dataset, which means that much of the required groundwork must be
constructed from scratch. The third, and arguably most critical, concerns reliability. A
language model permitted to answer independently may fabricate a rule that does not
exist. In the context of legal information, such an error is not a minor inconvenience;
it can mislead a person in a consequential situation.

\hs Taken together, these challenges indicate that an existing legal AI system cannot
simply be applied to the Cambodian context without modification. Both the retrieval and
the generation components must be adapted carefully to suit the Khmer language and the
particular structure of Cambodian law. Resolving this problem is the central concern of
this thesis.

% ─────────────────────────────────────────────────────────────────────
\newpage
\section{Objectives}

\hs The principal objective of this research is to design, implement, and evaluate an
AI-powered legal assistant for Cambodian law based on an Agentic Retrieval-Augmented
Generation architecture. The intention is not merely to describe such a system in
theory, but to construct a working implementation and to measure its performance
rigorously. The specific objectives are as follows:

\begin{itemize}
    \item Design an agentic RAG pipeline with intent classification, query rewriting,
    hybrid retrieval, grounding verification, and answer generation stages orchestrated
    through a directed state graph.

    \item Implement a hybrid retrieval system combining dense semantic search (FAISS)
    and sparse keyword search (BM25), fused through Reciprocal Rank Fusion (RRF), to
    improve article retrieval across diverse query types in Khmer.

    \item Develop a two-step answer generation approach that separates Khmer legal
    prose generation from structured citation extraction to maximise answer quality
    and citation accuracy.

    \item Build a full-stack web application with a React frontend featuring citation
    chips and an integrated PDF viewer that opens source articles at the exact
    referenced page.

    \item Evaluate the system using a manually constructed 30-question dataset spanning
    10 legal categories, measuring Retrieval Precision@5, Citation Accuracy, and
    Grounding Rate.
\end{itemize}

% ─────────────────────────────────────────────────────────────────────
\section{Scope and Limitation}

\hs It is important to state clearly what this project addresses and what it does not.
The work concentrates on applying Agentic RAG to legal question answering in Khmer, using
Cambodian labour law as the initial domain. Labour law was selected because it affects a
large segment of the population on a daily basis, which makes it a meaningful and
practical test case. Nevertheless, the system has genuine limitations, and acknowledging
them openly is necessary for an honest assessment. The principal limitations are as
follows:

\begin{itemize}
    \item \textbf{Domain Coverage}: The current implementation indexes labour-related
    laws as a proof of concept. Other domains of Cambodian law such as civil, criminal,
    and commercial law are not yet indexed, though the architecture supports extension.

    \item \textbf{No Case Law or Judicial Interpretation}: The system answers based
    solely on statutory law text. Court rulings, judicial interpretations, and legal
    amendments after the document compilation date are not included.

    \item \textbf{Language Limitation}: The system is designed for Khmer-language
    queries. English or mixed-language queries may produce inconsistent results.

    \item \textbf{Vocabulary Mismatch}: Highly specialised sub-domains exhibit lower
    retrieval accuracy due to vocabulary distance between natural queries and formal
    legal text, a known limitation of embedding-based retrieval on low-resource
    languages.

    \item \textbf{Streaming Latency}: The system uses simulated token streaming; the
    full agent pipeline completes before any text is returned to the user, resulting
    in an average response latency of approximately 20 seconds.
\end{itemize}

\hs These limitations do not diminish the central contribution of the work. For the most
part, they indicate directions in which the system can be extended, several of which are
revisited in the future work section of this thesis.

% ─────────────────────────────────────────────────────────────────────
\section{Planning of Project}

\hs As with any research project, the work proceeded in stages rather than all at once.
The activities were distributed across the four months of the internship and divided
into phases that overlapped where appropriate. The Gantt chart presented below
summarises how the available time was allocated, beginning with the initial literature
review and data preparation, continuing through system implementation, and concluding
with testing, evaluation, and the writing of this report in preparation for the defense.
In practice, some phases required more time than originally anticipated while others
required less, which is to be expected once development is under way.

\begin{table}[H]
\centering
\resizebox{\textwidth}{!}{%
\begin{tabular}{|p{4.5cm}|*{16}{>{\centering\arraybackslash}p{0.5cm}|}}
\hline
\textbf{Task} & \multicolumn{16}{c|}{\textbf{Week}} \\
\cline{2-17}
 & \multicolumn{4}{c|}{Month 1} & \multicolumn{4}{c|}{Month 2} & \multicolumn{4}{c|}{Month 3} & \multicolumn{4}{c|}{Month 4} \\
\cline{2-17}
 & 1 & 2 & 3 & 4 & 5 & 6 & 7 & 8 & 9 & 10 & 11 & 12 & 13 & 14 & 15 & 16 \\
\hline
Literature Review &
\progressbar{ } & \progressbar{ } & \progressbar{ } & & & & & & & & & & & & & \\
\hline
Data Collection \& Preprocessing &
& & \progressbar{ } & \progressbar{ } & \progressbar{ } & & & & & & & & & & & \\
\hline
System Architecture Design &
& & & \progressbar{ } & \progressbar{ } & \progressbar{ } & & & & & & & & & & \\
\hline
Backend Implementation &
& & & & \progressbar{ } & \progressbar{ } & \progressbar{ } & \progressbar{ } & \progressbar{ } & & & & & & & \\
\hline
Frontend Implementation &
& & & & & & \progressbar{ } & \progressbar{ } & \progressbar{ } & \progressbar{ } & & & & & & \\
\hline
System Integration \& Testing &
& & & & & & & & \progressbar{ } & \progressbar{ } & \progressbar{ } & & & & & \\
\hline
Evaluation \& Analysis &
& & & & & & & & & \progressbar{ } & \progressbar{ } & \progressbar{ } & \progressbar{ } & & & \\
\hline
Thesis Writing &
& \progressbar{ } & \progressbar{ } & \progressbar{ } & \progressbar{ } & \progressbar{ } &
\progressbar{ } & \progressbar{ } & \progressbar{ } & \progressbar{ } & \progressbar{ } &
\progressbar{ } & \progressbar{ } & \progressbar{ } & & \\
\hline
Defense Preparation &
& & & & & & & & & & & & \progressbar{ } & \progressbar{ } & \progressbar{ } & \progressbar{ } \\
\hline
\end{tabular}%
}
\caption{Project Roadmap}
\label{table:gantt}
\end{table}
